\section{The analytic inputs and the lifting problem}\label{sec:analytic-inputs}

The geometry will produce a divisor representing $K_Y$ with coefficients
of both signs. We shall prove that its full support is of logarithmic
general type. The difficulty is to find the first section of an auxiliary
nef adjoint before invoking a semiampleness criterion. The two analytic
results below, proved in the complete companion \cite{OpenAIMetrics2026},
address exactly that difficulty. They concern actual rational line bundles.

A semipositive singular Hermitian metric has \emph{minimal singularities}
if its local weight is, up to a bounded additive error, no more singular
than the weight of any other semipositive metric on the same bundle.
For a rational line bundle this definition is made after taking one
Cartier multiple and dividing the weights by that multiple.

\begin{theorem}[Zero Lelong numbers for nef klt adjoints]
\label{thm:adjoint-metric}
Let $(H,\Theta)$ be a normal connected projective complex klt pair,
where $\Theta\ge0$ is rational and $D_H=K_H+\Theta$ is nef and
$\Q$-Cartier. Let $\pi:V\to H$ be any projective log resolution.
The rational line bundle $\pi^*D_H$ admits a semipositive metric
with minimal singularities, and every such metric has zero Lelong
number at every point of $V$. Its multiplier ideal is trivial for
every positive exponent. Neither $K_H$ nor $\Theta$ is required
to be separately $\Q$-Cartier.
\end{theorem}

This is \cite[Theorem~1.1]{OpenAIMetrics2026}. The adjoint hypothesis is material:
a general nef line bundle need not have this metric regularity
\cite[Example~1.7]{DPS1994}.

\begin{theorem}[Ordinary interior injectivity]
\label{thm:interior-injectivity}
Let $V$ be a smooth projective complex variety and $L$ an integral
divisor. Let $0\le C_0\le C_2$ be effective rational divisors whose
combined support has simple normal crossings. Suppose the actual
rational line bundles $L-C_0$ and $L-C_2$ have semipositive singular
Hermitian metrics with zero Lelong numbers at every point. For a
rational number $0<\lambda<1$ put $C_1=(1-\lambda)C_0+\lambda C_2$.
The natural inclusion of sheaves induces an injection
\[
 H^1(V,\cO_V(K_V+L-\lfloor C_1\rfloor))
 \longrightarrow H^1(V,\cO_V(K_V+L-\lfloor C_0\rfloor)).
\]
\end{theorem}

The complete proof, including its passage to ordinary coherent
cohomology, is in \cite[Theorem~1.2]{OpenAIMetrics2026}. The analytic methods have antecedents
in Fujino's singular-metric treatment of harmonic forms and Matsumura's
uniform control of primitives \cite[Theorem~1.2 and Lemmas~3.1--3.2]
{FujinoInjectivity2012} and \cite[Theorems~1.3 and~5.9]{Matsumura2018}.
Those results are not being identified with the interior comparison.

Here is the lifting problem to which we will apply the theorem.
For a nef rational adjoint $N$ on a normal projective variety $H$
and a reduced closed subscheme $S\subset H$, choose $m>0$ with
$mN$ Cartier. The restriction sequence is
\[
 0\longrightarrow\cO_H(mN)\otimes\cI_S
 \longrightarrow\cO_H(mN)
 \longrightarrow\cO_S(mN)\longrightarrow0.
\]
The obstruction to lifting a section on $S$ is its image under the
connecting homomorphism, whose image is the kernel of
\[
 H^1(H,\cO_H(mN)\otimes\cI_S)
 \longrightarrow H^1(H,\cO_H(mN)).
\]
Thus injectivity of this ordinary cohomology map suffices to lift
every section on the whole scheme $S$. In the next section, $S$ is
the reduced non-klt floor of an auxiliary lc boundary. Two nearby
nef klt adjoints supply the endpoint metrics, and multiplier-ideal
local vanishing identifies the theorem's map with this inclusion.
Threefold abundance supplies a section on the entire floor, with
the conductor identifications already imposed.
